\documentclass[11pt]{article}
% =========================================================
%  學生講義 共用前言（以 XeLaTeX 編譯）
%  需要套件：xeCJK, tcolorbox（其餘多在 BasicTeX 內）
% =========================================================
\usepackage[a4paper,margin=2.3cm]{geometry}
\usepackage{amsmath,amssymb}
\usepackage{xcolor}
\usepackage{graphicx}

% ---- 中文字型（macOS 系統字型）----
\usepackage{xeCJK}
\setCJKmainfont{Songti TC}      % 內文：宋體
\setCJKsansfont{Heiti TC}       % 標題/強調：黑體
\xeCJKsetup{AutoFakeBold=2.2, AutoFakeSlant=0.2}

% ---- 版面 ----
\setlength{\parindent}{0pt}
\setlength{\parskip}{0.55em}
\linespread{1.15}
\renewcommand{\baselinestretch}{1.15}

% ---- 顏色 ----
\definecolor{cBlue}{HTML}{1f6feb}
\definecolor{cGray}{HTML}{6b7280}
\definecolor{cGrayBg}{HTML}{f3f4f6}
\definecolor{cPurp}{HTML}{7a3ff2}
\definecolor{cPurpBg}{HTML}{f5f1ff}

% ---- 標註框 ----
\usepackage[most]{tcolorbox}

% 就地補充新概念（灰底）：\begin{補充框}{標題}
\newtcolorbox{補充框}[1]{
  enhanced, breakable, arc=2pt, boxrule=0.6pt,
  colback=cGrayBg, colframe=cGray,
  left=9pt,right=9pt,top=7pt,bottom=7pt,
  before upper={\textbf{\sffamily #1}\par\vspace{3pt}},
}

% 延伸 / 開放問題（紫色左邊條）：\begin{延伸框}{標題}
\newtcolorbox{延伸框}[1]{
  enhanced, breakable, arc=2pt, boxrule=0pt,
  colback=cPurpBg, colframe=cPurpBg,
  borderline west={3pt}{0pt}{cPurp},
  left=10pt,right=9pt,top=7pt,bottom=7pt,
  before upper={\textbf{\sffamily 延伸閱讀｜#1}\par\vspace{3pt}},
}

% ---- 圖：置中、底下小字說明 ----
% \fig{檔名}{寬度}{說明}
\newcommand{\fig}[3]{%
  \begin{center}
  \includegraphics[width=#2\linewidth]{#1}\\[2pt]
  {\small\color{cGray} #3}
  \end{center}}

% ---- 章節標題用黑體 ----
\newcommand{\h}[1]{\sffamily #1}


\begin{document}

\begin{center}
{\sffamily\Huge\bfseries 必物-3 物質的組成與交互作用}\\[3pt]
{\sffamily\large 普通高中 物理（必修）・學生講義}
\end{center}
\vspace{0.6em}

本章把一塊看得見、摸得到的物質\textbf{一層一層拆開}：分子、原子、原子核與電子、質子與中子，最後到\textbf{夸克（quark）}。每拆開一層，都要回答同一個問題——\textbf{是什麼力把這一層綁在一起}。把這些力歸納起來，自然界的各種作用最後只歸結為\textbf{四種基本交互作用（four fundamental interactions）}：強力、電磁力、弱力、重力。

前一章學到的接觸力（正向力、摩擦力、張力）其實都是電磁力的不同表現；本章正式回答「力的微觀來源」。其中電子靠\textbf{電力}被綁在核外，原子核則靠\textbf{強力}維持穩定。本章重觀念、輕計算，內容以\textbf{定性理解}為主。

\medskip
本章主線：

\begin{center}
物質的組成階層 $\rightarrow$ 原子結構（金箔散射）$\rightarrow$ 原子核與強力 $\rightarrow$ 四種基本交互作用
\end{center}

\section*{\sffamily 3-1\quad 物質的組成階層：從一塊冰到夸克}

\subsection*{\sffamily A.\ 物質是一層套一層的階層結構}

如果把一塊物質一直分割下去，會得到什麼？今天的答案是：物質具有\textbf{階層結構（hierarchical structure）}，一層套著一層。由大到小依序為：

$$\text{巨觀物質}\to\text{分子}\to\text{原子}\to
\begin{cases}\text{原子核}\to\text{質子 ＋ 中子}\to\text{夸克}\\[2pt]\text{電子（目前不可再分）}\end{cases}$$

\fig{m3-hierarchy}{0.72}{圖 3-1：物質的組成階層（由大到小）。每一層的右側標出該層的典型尺度；原子約 $10^{-10}$ m、原子核約 $10^{-15}$ m。}

有兩個尺度需要記住：

$$\boxed{\;\text{原子}\sim 10^{-10}\ \text{m}\,,\qquad \text{原子核}\sim 10^{-15}\ \text{m}\;}$$

兩者相差約 $10^{5}$（十萬）倍。原子核雖然只佔原子直徑的約十萬分之一，卻集中了原子\textbf{幾乎全部的質量}（質子、中子各約為電子質量的 1800 倍）。換言之，原子內部絕大部分是電子稀疏出沒的空間。可類比為：若把一個原子放大成一座體育場，原子核大約只有場中央一顆綠豆的大小。

\begin{補充框}{先說明：夸克是什麼？為什麼質子還能再分？}
質子與中子並非最基本的粒子，它們各由\textbf{三顆夸克}組成。夸克是目前已知\textbf{不可再分}的基本粒子之一。
\begin{itemize}\setlength{\itemsep}{1pt}
\item 質子 ＝ 兩顆\textbf{上夸克（up，$u$）}＋ 一顆\textbf{下夸克（down，$d$）}，記作 $uud$。
\item 中子 ＝ 一顆上夸克 ＋ 兩顆下夸克，記作 $udd$。
\end{itemize}
夸克帶\textbf{分數電荷}（$u$ 帶 $+\tfrac23 e$、$d$ 帶 $-\tfrac13 e$），組合後才湊出質子的 $+e$ 與中子的 $0$：

\quad 質子 $uud$：$\tfrac23+\tfrac23-\tfrac13=+1$（單位 $e$）；\qquad 中子 $udd$：$\tfrac23-\tfrac13-\tfrac13=0$。

\textbf{注意：}雖然質子由夸克組成，但目前無法把質子敲開、取得單獨一顆夸克。這並不矛盾——夸克之間的吸引力\textbf{不會隨距離拉遠而變弱}，可類比為一條拉不斷的橡皮筋，拉得越開、儲存的能量越多。當能量大到一定程度，依 $E=mc^2$，這份能量會直接轉變為一對新的夸克補在斷口，結果是得到新的粒子，而非單獨一顆夸克。因此夸克永遠被「關」在粒子內部，稱為\textbf{夸克禁閉（quark confinement）}。可類比為把一塊磁鐵分成兩半，永遠得到兩塊各有南北極的小磁鐵，無法分出單獨的磁極。夸克禁閉雖然反直覺，但已是\textbf{有定論}的現象：以高能電子探測質子內部，顯示質子內有三個點狀的散射中心，與夸克模型相符。
\end{補充框}

\textbf{注意：}「電子像行星繞太陽」只是早期的簡化圖像，並不精確。電子並非走確定的軌道，而是以機率分布的形式出現於核外空間（這點以及「電子為何不墜入核」的完整解釋見必物-6）。本章只用此圖像定性說明「電力把電子綁在核外」。另外，原子「幾乎是空的」並不表示物體可以互相穿透；阻擋物體互相穿過的是電子之間的\textbf{電磁排斥}，而非物質是否「實心」。

\section*{\sffamily 3-2\quad 原子的結構與金箔散射}

原子內部無法直接觀察，物理學家只能先\textbf{以模型推測，再用實驗檢驗}。原子模型的演進，是科學如何依據新證據修正自己的範例。

\subsection*{\sffamily A.\ 原子模型的演進}

\begin{center}
\renewcommand{\arraystretch}{1.35}
\begin{tabular}{p{0.10\linewidth} p{0.16\linewidth} p{0.32\linewidth} p{0.30\linewidth}}
\hline
\textbf{年代} & \textbf{提出者} & \textbf{模型圖像} & \textbf{被什麼修正} \\
\hline
1803 & 道耳吞（Dalton） & 原子是實心小球，不可再分 & 發現電子，原子有內部結構 \\
1897 起 & 湯姆森（Thomson） & 正電均勻分布，電子嵌於其中 & 金箔散射顯示正電高度集中 \\
1911 & 拉塞福（Rutherford） & 核式原子：正電與質量集中於極小的核 & 無法解釋原子穩定與光譜 \\
1913 & 波耳（Bohr） & 電子只能處於特定能階的軌道 & 由量子力學進一步取代 \\
\hline
\end{tabular}
\end{center}

道耳吞依化學定律（定比、倍比定律）論證原子存在；湯姆森發現了\textbf{電子（electron）}（帶負電、質量極小），於是原子不再被視為實心球，他推測正電均勻分布、電子嵌於其中。真正的轉折來自其學生拉塞福的一項實驗。

\subsection*{\sffamily B.\ 拉塞福金箔散射}

拉塞福以 \textbf{$\alpha$ 粒子}（帶正電、質量不小的高速粒子）轟擊極薄的金箔，觀察其被偏折的角度。實驗有三個關鍵觀測：

\begin{itemize}\setlength{\itemsep}{1pt}
\item \textbf{絕大多數} $\alpha$ 粒子幾乎直線穿過，方向幾乎不變。
\item \textbf{少數}被偏折較大角度。
\item \textbf{極少數（約萬分之一）}被反彈回來（偏折超過 $90^\circ$）。
\end{itemize}

\fig{m3-rutherford}{0.78}{圖 3-2：拉塞福金箔散射。多數 $\alpha$ 粒子直線穿過電子稀疏的空間；少數因接近原子核而大角度偏折，極少數正面接近原子核而被反彈。紅點代表原子核：極小、極重、帶正電，集中了幾乎全部質量。}

若正電如湯姆森模型般均勻、稀疏地分布，$\alpha$ 粒子最多只會被輕微偏折，\textbf{不可能被反彈}。能反彈，唯一的解釋是：原子的\textbf{正電與幾乎全部質量集中在一個極小、極重的核}裡。多數 $\alpha$ 粒子穿過的是核外電子稀疏的空間，只有正面接近這顆小核的才會被強烈的電斥力大幅偏折。由此得到\textbf{核式原子模型}：

$$\boxed{\;\text{原子}=\text{中央極小的正電核}\;+\;\text{核外的電子}\;}$$

原子核（nucleus）帶正電、集中幾乎全部質量，半徑約為原子的 $10^{-5}$ 倍；電子在核外廣大的空間運動。

\textbf{注意：}$\alpha$ 粒子被反彈，並非撞到電子。電子質量極小，可類比為以保齡球撞乒乓球，不會使保齡球反彈；使 $\alpha$ 粒子反彈的是又小又重又帶正電的原子核。此外，原子模型一個接一個被修正，並不代表科學結論不可靠：每個模型在其證據範圍內都成立，新證據出現時才被更精確的模型取代，這正是科學自我修正的過程。

\subsection*{\sffamily C.\ 維繫原子的力：電力}

原子核帶正電、電子帶負電，\textbf{異號電荷相吸}。正是這個\textbf{電力（庫侖吸引）}把電子留在核的周圍，形成穩定的原子。兩個帶電體間的電力可用\textbf{庫侖定律（Coulomb's law）}描述：與電荷量乘積成正比、與距離平方成反比，

$$F=k\,\frac{q_1 q_2}{r^2},$$

同號電荷相斥、異號相吸。原子中核（$+$）與電子（$-$）異號，因此相吸，電子被束縛於核外。

電力提供電子被核束縛所需的吸引，這點與重力把行星留在太陽旁類似。差別在於：在原子尺度下，電力遠強於重力。以氫原子為例，質子與電子之間的電力約為兩者間重力的 $10^{39}$ 倍，因此討論原子結構時可完全忽略重力。這也預告了 3-4 的主題：四種基本力的強弱差距極大。

\section*{\sffamily 3-3\quad 原子核與強力}

\subsection*{\sffamily A.\ 一個矛盾：核為什麼不被電斥力推開}

原子核內擠著多個質子（氦有 2 個、碳有 6 個、鈾有 92 個），它們\textbf{全帶正電}。同號電荷相斥，且核內距離極近（$\sim10^{-15}$ m），由 $F\propto 1/r^2$，這麼近的電斥力相當大。僅就電力而言，原子核應被自己的電斥力推開。

但大多數原子核長期穩定。這表示\textbf{核內必有一種更強的吸引力，在極近距離把核子們拉住}，這就是\textbf{強力（strong force，又稱強核力）}。

\fig{m3-nuclear-balance}{0.72}{圖 3-3：原子核內電斥力與強力的平衡。質子間的電力互相排斥（往外）；質子、中子間的強力互相吸引（往內）。核內的近距離下強力勝過電斥力，原子核因此穩定。中子不帶電、只感受強力，協助維繫原子核。}

強力有兩個特徵：
\begin{itemize}\setlength{\itemsep}{1pt}
\item \textbf{極強}：在原子核的尺度上，強到足以壓過質子間的電斥力，把核子拉在一起。
\item \textbf{極短程}：作用範圍只有約 $10^{-15}$ m（一個核子的大小）。超出此距離，強力幾乎立即趨於零，衰減比電力的 $1/r^2$ 快得多。
\end{itemize}

$$\boxed{\;\text{近距離（核內）：強力 > 電斥力}\ \Rightarrow\ \text{原子核穩定}\;}$$

「又強又短程」是關鍵：強力只在緊鄰的核子間有效，無法像電力那樣作用到遠處；但在核內的近距離，它穩定地勝過電斥力。可類比為強力像必須貼近才有效的強力磁鐵，電力則像作用較弱但作用範圍較遠的彈簧。

\subsection*{\sffamily B.\ 中子的角色}

中子不帶電，\textbf{不參與電斥力}，卻同樣感受強力。增加中子等於\textbf{只增加吸引、不增加排斥}，有助於把核束縛得更穩。這說明為什麼較大的原子核需要\textbf{更多中子}：質子越多，須對抗的電斥力越大，因此需要更多中子提供額外的強力。

\begin{補充框}{延伸說明：為什麼有些原子核仍會不穩定（衰變、分裂）？}
強力雖強，但\textbf{短程}。當質子數很多、原子核很大時，外圍的質子可能已超出強力的有效範圍，但\textbf{長程的電斥力仍持續累積}。當質子多到一定程度（如鈾有 92 個），電斥力的長程累積終於勝過短程的強力，原子核就變得不穩定。這正是放射性衰變與核分裂的根源。其中能量的釋放牽涉質量虧損與 $E=mc^2$，屬於必物-5 的主題；本章只說明強力維繫原子核穩定這一面。
\end{補充框}

\subsection*{\sffamily C.\ 質子數、中子數與同位素}

\begin{itemize}\setlength{\itemsep}{1pt}
\item \textbf{質子數 $Z$（原子序）}：決定元素種類（氫 $Z=1$、氦 $Z=2$、碳 $Z=6$$\dots$）。
\item \textbf{中子數 $N$}：同一元素可有不同的中子數。
\item \textbf{質量數 $A=Z+N$}：質子數加中子數。
\item \textbf{同位素（isotope）}：質子數\textbf{相同}、中子數\textbf{不同}的原子，是同一元素的不同版本。其化學性質幾乎相同（化學性質由電子決定，而電子數等於質子數），但質量不同、原子核的穩定性也可能不同。
\end{itemize}

記法為 ${}^{A}_{Z}\mathrm{X}$，例如 ${}^{12}_{6}\mathrm{C}$（碳-12）與 ${}^{14}_{6}\mathrm{C}$（碳-14）是碳的兩種同位素：質子數同為 6，中子數分別為 6 與 8。

\textbf{注意：}強力只在原子核尺度（$\sim10^{-15}$ m）內作用，日常生活中的「黏」「撐」「拉」都屬於電磁力，與強力無關。此外，中子之間、中子與質子之間都有強力，並非「中子不帶電就沒有任何作用力」；中子只是不參與電力。

\section*{\sffamily 3-4\quad 自然界的四種基本交互作用}

把前面三節的力收攏起來看：維繫電子的是電力、維繫原子核的是強力、把物體拉向地面的是重力。自然界的各種作用最後歸結為\textbf{四種基本交互作用}。

\begin{center}
\renewcommand{\arraystretch}{1.35}
\begin{tabular}{p{0.14\linewidth} p{0.20\linewidth} p{0.20\linewidth} p{0.30\linewidth}}
\hline
\textbf{交互作用} & \textbf{相對強度（量級）} & \textbf{作用範圍} & \textbf{負責什麼} \\
\hline
\textbf{強力} & $1$（最強） & 極短，$\sim10^{-15}$ m & 把質子、中子綁成原子核；把夸克綁成核子 \\
\textbf{電磁力} & $\sim10^{-2}$ & 無限遠（$\propto 1/r^2$） & 維繫原子、化學鍵；日常接觸力、摩擦、彈力、光 \\
\textbf{弱力} & $\sim10^{-13}$ & 極短，$\sim10^{-18}$ m & 主導某些放射性衰變（如 $\beta$ 衰變）、恆星核反應的關鍵步驟 \\
\textbf{重力} & $\sim10^{-38}$（最弱） & 無限遠（$\propto 1/r^2$） & 主宰天體與宇宙的結構 \\
\hline
\end{tabular}
\end{center}

\fig{m3-forces}{0.82}{圖 3-4：四種基本交互作用的相對強度（以強力為 $1$，橫軸為 10 的次方）。四者強度差距極大，從強力的 $1$ 到重力的約 $10^{-38}$。圖中同時標出各力的作用範圍。}

相對強度為\textbf{數量級的粗略比較}，數值隨比較方式略有不同，重點是\textbf{彼此差距極大}。

\subsection*{\sffamily A.\ 「強弱」與「作用範圍」要分開看}

一個力\textbf{強不強}和它\textbf{作用得多遠}是兩件不同的事：
\begin{itemize}\setlength{\itemsep}{1pt}
\item 強力與弱力都是\textbf{短程}的：影響力出了原子核就幾乎消失，因此日常生活感覺不到。
\item 電磁力與重力都是\textbf{長程}的（$\propto1/r^2$，距離再遠都不會真正歸零），能跨越遙遠距離持續累積。
\end{itemize}

\subsection*{\sffamily B.\ 為什麼最弱的重力卻主宰宇宙}

重力是四種力中最弱的（比電磁力弱約 $10^{36}$ 倍），但在天文尺度上由它主導。原因有兩個：

\textbf{第一，重力不會抵消。}質量只有一種，重力\textbf{永遠是吸引、只會累加}。電荷則有正負兩種，大塊物質（行星、恆星）整體接近電中性，正負電荷的作用大半互相抵消，對外幾乎不顯電磁力。可類比為拔河：電磁力像方向不一致的一群人，合力幾乎為零；重力像所有人都往同一方向拉，即使每人力量小，人數一多合力仍很大。

\textbf{第二，重力是長程的。}它不像強力、弱力出了原子核就失效，而能一路累積到星系尺度。

於是：\textbf{微觀世界由電磁力與強力主宰，宏觀宇宙由重力主宰}。重力主宰宇宙並非因為它強，而是因為它不抵消又作用得遠。

\textbf{注意：}「重力最弱所以不重要」並不正確。在天文尺度重力是主導因素。「強弱」指的是\textbf{同樣兩個粒子、同樣距離下}單對粒子的比較；「能否主宰某尺度」則要看累積後的總效果，須把「會不會抵消」與「作用得多遠」一併考慮。另外，磁力並非獨立的第五種力——電與磁是同一種交互作用的兩面，合稱\textbf{電磁力}（這是必物-4 的主題）；弱力也不是「較弱的電磁力」，而是性質不同的另一種基本力，負責某些衰變。

\begin{延伸框}{四種基本力能不能統一成一種？}
物理學追求\textbf{用更少的原理解釋更多的現象}。「統一」即發現「兩種看似不同的力其實是同一種力在不同條件下的表現」。這條路已走了數步，但只完成一部分：

\begin{center}
\renewcommand{\arraystretch}{1.3}
\begin{tabular}{p{0.22\linewidth} p{0.40\linewidth} p{0.24\linewidth}}
\hline
\textbf{統一} & \textbf{把哪些合而為一} & \textbf{現況} \\
\hline
電磁統一 & 電力 ＋ 磁力 $\rightarrow$ 電磁力 & 已完成、定論 \\
電弱統一 & 電磁力 ＋ 弱力 $\rightarrow$ 電弱作用 & 已完成、定論 \\
大統一（GUT） & 電弱 ＋ 強力 $\rightarrow$ 一種力 & 尚無定論 \\
萬有理論（TOE） & 再加上重力 $\rightarrow$ 一種力 & 尚無定論 \\
\hline
\end{tabular}
\end{center}

\textbf{已是定論的部分：}馬克士威方程式把電、磁、光收進同一套理論，完成\textbf{電磁統一}。1960 年代的理論預言，在足夠高的能量下電磁力與弱力其實是同一種「電弱作用」的兩面；其後相關粒子（W、Z 玻色子）與希格斯玻色子陸續被實驗發現，\textbf{電弱統一已成定論}。至此四種力已併成三種（電弱、強力、重力）。

\textbf{仍是開放問題的部分：}把強力也納入的「大統一理論（GUT）」有不少候選模型，但\textbf{至今沒有一個被實驗證實}；它預言的某些現象（如質子衰變）一直未被觀測到。把重力也納入的「萬有理論（Theory of Everything）」更困難：描述重力的廣義相對論與描述其他三力的量子理論，在數學架構上\textbf{至今無法相容}。弦論、迴圈量子重力等是候選方向，但都尚無實驗支持。

簡言之：四種力已從四種收斂到三種，但能否再收到一種，\textbf{目前尚無定論}。這是當代理論物理最前沿的開放問題之一。
\end{延伸框}

\section*{\sffamily 3-5\quad 本章重點整理}

\begin{補充框}{一頁複習（關鍵結論）}
\begin{itemize}\setlength{\itemsep}{2pt}
\item \textbf{組成階層}：巨觀物質 $\rightarrow$ 分子 $\rightarrow$ 原子 $\rightarrow$ 原子核＋電子 $\rightarrow$ 質子＋中子 $\rightarrow$ 夸克。兩個必記尺度：原子 $\sim10^{-10}$ m、原子核 $\sim10^{-15}$ m（差 $10^5$ 倍，原子內部大部分是空的）。
\item \textbf{原子模型}：道耳吞（實心球）$\rightarrow$ 湯姆森（正電均勻分布）$\rightarrow$ 拉塞福（核式，由金箔散射確立）$\rightarrow$ 波耳（能階）。
\item \textbf{金箔散射三觀測}：多數直穿、少數大偏折、極少數反彈，顯示正電與質量集中於極小的原子核。
\item \textbf{維繫原子的力＝電力}：核（$+$）吸引電子（$-$），$F=kq_1q_2/r^2$；同尺度下電力遠強於重力（氫原子裡約 $10^{39}$ 倍）。
\item \textbf{維繫原子核的力＝強力}：又強又短程（$\sim10^{-15}$ m），近距離勝過質子間電斥力，原子核因而穩定；中子只感受強力、不帶電，協助維繫原子核。
\item \textbf{同位素}：$Z$ 相同、$N$ 不同；$A=Z+N$；$Z$ 決定元素種類。
\item \textbf{四種基本交互作用}：強力（$1$，極短）> 電磁力（$10^{-2}$，長程）> 弱力（$10^{-13}$，極短）> 重力（$10^{-38}$，長程）。微觀靠電磁力與強力，宇宙靠重力（因重力不抵消又長程）。
\item \textbf{統一}：電磁統一、電弱統一已是定論；大統一與含重力的萬有理論目前尚無定論。
\end{itemize}
\end{補充框}

\end{document}
